\begin{studiumworked}{M18. Gas inerte en expansión: dos límites térmicos · INTERMEDIATE}
\textbf{Enunciado original.} Una masa fija de gas ideal inerte ocupa inicialmente un volumen $V_0$, a 2 MPa absolutos y 300 K. El volumen aumenta a $1.8V_0$. Halle presión y temperatura para expansión adiabática reversible con gamma 1.4 y compare presión isotérmica.

\textbf{Datos utilizados.} Expansión: $1.8$; Presión absoluta: $2.000000\times10^{6}\,\mathrm{Pa}$; Temperatura inicial: $300\,\mathrm{K}$.

\textbf{Incógnitas.} Magnitudes solicitadas y cierre físico.

\textbf{Hipótesis.} Modelo académico con hipótesis del capítulo. Datos originales; no representan hardware ni ensayo.

\textbf{Modelo físico.} Conservación y relaciones constitutivas declaradas.
\[p/p_0=(V/V_0)^{-\gamma},\quad T/T_0=(V/V_0)^{1-\gamma}\]
\[p_{iso}/p_0=V_0/V\]
\textbf{Despejes y sustituciones.}
\begin{enumerate}
\item El proceso adiabático da $p=2\times10^6(1.8)^{-1.4}=8.783113\times10^{5}$ Pa y $T=300(1.8)^{-0.4}=237.1441$ K.
\item El isotérmico conserva 300 K y da $p=2\times10^6/1.8=1.111111\times10^{6}$ Pa. Su presión mayor requiere entrada de calor para compensar el trabajo de expansión.
\item La masa se mantiene fija en ambas hipótesis; este cálculo no debe trasladarse sin cambios a un reservorio rígido que descarga gas.
\end{enumerate}
\textbf{Resultados.} $p_{ad}=8.783113\times10^{5}$ Pa; $T_{ad}=237.1441$ K; $p_{iso}=1.111111\times10^{6}$ Pa.

\textbf{Interpretación.} El resultado corresponde a las hipótesis indicadas; su interpretación requiere conservar fronteras y unidades.

\textbf{Comprobación.} En cada proceso, pV/(RT) recupera la misma masa inicial; ambos estados tienden al inicial cuando la razón de volumen tiende a uno.

\textbf{Límites.} Gas ideal uniforme; sin líquido, disolución ni red de alimentación.
\end{studiumworked}
